\section{Presentación del invitado y contexto de la entrevista}

El invitado de esta edición del WhynotTV Podcast (la tercera) es Chen Tianqi (陈天奇, Tianqi Chen), profesor asistente del departamento de aprendizaje automático de CMU y cofundador de OctoML (adquirida por NVIDIA a fines de 2024). Durante casi 20 años, Chen Tianqi ha impulsado, en torno al eje de \textbf{``hacer el aprendizaje automático más ligero, más rápido y más fácil de desplegar''}, proyectos de código abierto ampliamente adoptados como XGBoost, MXNet, TVM y MLC LLM. El artículo de XGBoost, en particular, tiene cerca de 70,000 citas y sigue siendo la herramienta más usada en el ámbito de los datos tabulares.

En 2019 publicó en Zhihu el artículo «Diez años de investigación en aprendizaje automático», que recibió decenas de miles de «me gusta». Seis años después, en 2025, ya lleva casi 20 años en este camino del aprendizaje automático. El conductor (WhynotTV) señala: ``Para mí, Chen Tianqi no es solo un nombre reconocido en el ámbito de la IA; también marcó mi primera comprensión del aprendizaje automático y de la investigación científica.''

Esta entrevista dura 160 minutos y parte de su infancia, pasa por la clase ACM de licenciatura, el laboratorio APEX de posgrado, los cuatro grandes proyectos de su doctorado en UW (XGBoost, MXNet, TVM y Net2Net durante su estancia en Google Brain), hasta llegar a su puesto docente en CMU y su emprendimiento con OctoML, para finalmente elevarse a una filosofía de vida sobre el largo plazo, la intención original y el valor.

\begin{importantbox}{El eje técnico de Chen Tianqi}
De XGBoost (2014) a MXNet (2015), a TVM (2017) y luego a MLC LLM (2023), cada proyecto de Chen Tianqi gira en torno al mismo problema central: \textbf{cómo lograr que los sistemas de aprendizaje automático funcionen de manera eficiente en distintos tipos de hardware, con menos recursos de ingeniería y de forma más automatizada}. Este eje atraviesa la era del big data, la era del aprendizaje profundo y la era de los modelos grandes. No es un investigador atado a un método en particular, sino alguien que persigue un problema------los métodos pueden cambiar (de los modelos de árboles a los marcos de aprendizaje profundo y a los compiladores), pero el problema sigue siendo el mismo: ``mejorar los sistemas de ML''.
\end{importantbox}

